\documentclass[../../main]{subfiles}

\begin{document}

\section{Fechos de Relações}
\label{sec:matematica:relacoes:fechos-de-relacoes}

\begin{defn}[Fecho]
	\label{def:matematica:relacoes:fechos:fecho}

	Seja \(R\) uma relação binária sobre um conjunto \(A\).

	Seja \(P\) uma propriedade de relações.

	O fecho de \(R\) com respeito à propriedade \(P\) é a menor relação
	\(S\) sobre \(A\) tal que

	\[
		R \subseteq S
	\]

	e \(S\) satisfaz a propriedade \(P\).

\end{defn}

\begin{obs}

	A expressão ``menor relação'' significa que, para qualquer relação
	\(T\) que contenha \(R\) e satisfaça a propriedade considerada,

	\[
		S \subseteq T.
	\]

\end{obs}

\subsection{Fecho Reflexivo}

\begin{defn}[Fecho Reflexivo]
	\label{def:matematica:relacoes:fechos:fecho-reflexivo}

	Seja \(R\) uma relação sobre \(A\).

	O fecho reflexivo de \(R\) é a relação

	\[
		R_r
		=
		R
		\cup
		\{
		(a,a)
		\mid
		a \in A
		\}.
	\]

\end{defn}

\begin{prop}
	\label{prop:matematica:relacoes:fechos:fecho-reflexivo-reflexivo}

	O fecho reflexivo de uma relação é reflexivo.

\end{prop}

\begin{proof}

	Seja \(a \in A\).

	Pela definição,

	\[
		(a,a)
		\in
		\{
		(x,x)
		\mid
		x \in A
		\}.
	\]

	Portanto,

	\[
		(a,a)\in R_r.
	\]

	Logo \(R_r\) é reflexiva.

\end{proof}

\begin{prop}
	\label{prop:matematica:relacoes:fechos:minimalidade-fecho-reflexivo}

	O fecho reflexivo é a menor relação reflexiva que contém \(R\).

\end{prop}

\begin{proof}

	Seja \(T\) uma relação reflexiva tal que

	\[
		R \subseteq T.
	\]

	Como \(T\) é reflexiva,

	\[
		(a,a)\in T
	\]

	para todo \(a\in A\).

	Portanto,

	\[
		R_r \subseteq T.
	\]

\end{proof}

\subsection{Fecho Simétrico}

\begin{defn}[Relação Inversa]
	\label{def:matematica:relacoes:fechos:relacao-inversa}

	Seja \(R\) uma relação sobre \(A\).

	Define-se a relação inversa de \(R\) por

	\[
		R^{-1}
		=
		\{
		(b,a)
		\mid
		(a,b)\in R
		\}.
	\]

\end{defn}

\begin{defn}[Fecho Simétrico]
	\label{def:matematica:relacoes:fechos:fecho-simetrico}

	O fecho simétrico de \(R\) é a relação

	\[
		R_s
		=
		R
		\cup
		R^{-1}.
	\]

\end{defn}

\begin{prop}
	\label{prop:matematica:relacoes:fechos:fecho-simetrico-simetrico}

	O fecho simétrico de uma relação é simétrico.

\end{prop}

\begin{proof}

	Seja

	\[
		(a,b)\in R_s.
	\]

	Então

	\[
		(a,b)\in R
	\]

	ou

	\[
		(a,b)\in R^{-1}.
	\]

	Em ambos os casos,

	\[
		(b,a)\in R_s.
	\]

	Portanto \(R_s\) é simétrica.

\end{proof}

\begin{prop}
	\label{prop:matematica:relacoes:fechos:minimalidade-fecho-simetrico}

	O fecho simétrico é a menor relação simétrica que contém \(R\).

\end{prop}

\begin{proof}

	Seja \(T\) uma relação simétrica contendo \(R\).

	Como \(T\) é simétrica,

	\[
		R^{-1}\subseteq T.
	\]

	Portanto,

	\[
		R_s
		=
		R\cup R^{-1}
		\subseteq T.
	\]

\end{proof}

\subsection{Fecho Transitivo}

\begin{defn}[Potência de uma Relação]
	\label{def:matematica:relacoes:fechos:potencia-relacao}

	Seja \(R\) uma relação sobre \(A\).

	Define-se

	\[
		R^1 = R
	\]

	e, para \(n\ge 1\),

	\[
		R^{n+1}
		=
		R^n \circ R,
	\]

	onde \(\circ\) denota a composição de relações.

\end{defn}

\begin{defn}[Fecho Transitivo]
	\label{def:matematica:relacoes:fechos:fecho-transitivo}

	O fecho transitivo de \(R\) é a relação

	\[
		R_t
		=
		\bigcup_{n=1}^{\infty} R^n.
	\]

\end{defn}

\begin{obs}

	Um par

	\[
		(a,b)\in R_t
	\]

	quando existe um caminho finito ligando \(a\) a \(b\) por elementos de
	\(R\).

\end{obs}

\begin{prop}
	\label{prop:matematica:relacoes:fechos:fecho-transitivo-transitivo}

	O fecho transitivo de uma relação é transitivo.

\end{prop}

\begin{proof}

	Se

	\[
		(a,b)\in R^m
	\]

	e

	\[
		(b,c)\in R^n,
	\]

	então

	\[
		(a,c)\in R^{m+n}.
	\]

	Como

	\[
		R^{m+n}
		\subseteq
		R_t,
	\]

	segue que

	\[
		(a,c)\in R_t.
	\]

	Portanto \(R_t\) é transitiva.

\end{proof}

\begin{prop}
	\label{prop:matematica:relacoes:fechos:minimalidade-fecho-transitivo}

	O fecho transitivo é a menor relação transitiva que contém \(R\).

\end{prop}

\begin{proof}

	Seja \(T\) uma relação transitiva contendo \(R\).

	Por indução sobre \(n\), obtém-se

	\[
		R^n \subseteq T
	\]

	para todo \(n\ge 1\).

	Portanto,

	\[
		R_t
		=
		\bigcup_{n=1}^{\infty}R^n
		\subseteq T.
	\]

\end{proof}

\begin{ex}

	Seja

	\[
		A=\{1,2,3\}
	\]

	e

	\[
		R
		=
		\{
		(1,2),
		(2,3)
		\}.
	\]

	Como

	\[
		(1,2)\in R
		\quad\text{e}\quad
		(2,3)\in R,
	\]

	o fecho transitivo deve conter

	\[
		(1,3).
	\]

	Assim,

	\[
		R_t
		=
		\{
		(1,2),
		(2,3),
		(1,3)
		\}.
	\]

\end{ex}

\end{document}
